\section{The critical exponent: small and large moment budgets}
At $s=s_c=(1+\rho)/(1-\rho)$ the size of the moment budget matters.
We give a sufficient stability condition near $B=1$ and a counterexample
for every $B>2$. These leave intermediate budgets unresolved; no optimal
critical budget is asserted. All norms in this section use the uniform
probability measure $u_L$.

\begin{theorem}[Small critical budgets]\label{thm:critical-small}
Fix $0<\rho<1$ and $s=s_c$. For any $L\ge1$ and any density $r\ge0$
with $\mathbb E r=1$ and $\|r\|_s\le B$, $B\ge1$,
\[
 \gap(T_r)\ge\max\{0,1-\rho-d_s(B)\},\qquad
 d_s(B)=\begin{cases}
 (B^s-1)^{1/s},&s\ge2,\\
 \sqrt{\dfrac{2(B^s-1)}{s(s-1)}}(B+1)^{(2-s)/2},&1<s<2.
 \end{cases}
\]
In particular, each fixed $\rho$ admits a nontrivial interval of budgets
$B>1$ with a positive dimension-independent gap. No pointwise density
bounds are required.
\end{theorem}
\begin{proof}
Put $Q=P_{\sqrt\rho}$ and $S_r=QM_rQ$, which has the spectrum of $T_r$.
Ordinary Boolean hypercontractivity, recalled in~\cite[Eq.~(1.1)]{MOS},
gives $\|Qf\|_q\le\|f\|_2$ for $q=2s/(s-1)$, since
$\rho=(s-1)/(s+1)=1/(q-1)$. H\"older's inequality with exponents
$s,q,q$ yields
\[
 |\langle f,Q M_{r-1}Qg\rangle|
 \le\|r-1\|_s\|Qf\|_q\|Qg\|_q
 \le\|r-1\|_s\|f\|_2\|g\|_2.
\]
Thus $\|S_r-P_\rho\|\le\|r-1\|_s$ and
$\lambda_2(S_r)\le\rho+\|r-1\|_s$. Also
$\lambda_1(S_r)\ge\langle1,S_r1\rangle=1$.
It remains to show $\|r-1\|_s\le d_s(B)$.

Define $\phi_s(t)=t^s-1-s(t-1)$ for $t\ge0$, so that
$\mathbb E\phi_s(r)\le B^s-1$. For $s\ge2$ and $u\ge0$,
\[
 \phi_s(1+u)=s\int_0^u((1+v)^{s-1}-1)\,dv\ge u^s,
\]
using $(1+v)^{s-1}\ge1+v^{s-1}$. For $0\le u\le1$,
the derivative of $\phi_s(1-u)-u^s$ is
$s[1-(1-u)^{s-1}-u^{s-1}]\ge0$ and its value at zero is zero.
Hence $\phi_s(t)\ge|t-1|^s$, proving the first estimate by integration.
For $1<s<2$, the integral Taylor remainder gives
\begin{align*}
 \phi_s(t)&=s(s-1)(t-1)^2
  \int_0^1(1-v)[1+v(t-1)]^{s-2}\,dv\\
 &\ge\frac{s(s-1)}2\frac{(t-1)^2}{(t+1)^{2-s}}.
\end{align*}
At $t=0$ the identity follows by continuity; the integral is finite for
$s>1$. For $W=(r-1)^2/(r+1)^{2-s}$, H\"older gives
\[
 \mathbb E|r-1|^s
 \le(\mathbb EW)^{s/2}
       (\mathbb E(r+1)^s)^{(2-s)/2}.
\]
Use $\mathbb EW\le2(B^s-1)/(s(s-1))$ and $\|r+1\|_s\le B+1$.
Both formulas for $d_s(B)$ tend to zero as $B\downarrow1$.
Positive semidefiniteness supplies the trivial zero lower bound.
\end{proof}

\begin{theorem}[Balanced critical tilt]\label{thm:critical-large}
Fix $0<\rho<1$, $s=s_c$, and $B>2$. There are constants $c,d>0$,
independent of $L$, and strictly positive densities $r_L$, for every $L\ge1$,
such that
\[
 \mathbb E r_L=1,\quad\|r_L\|_s\le B,\quad r_L\ge1/B,
 \qquad
 0<\gap(T_{r_L})\le2\sqrt{B-1}\,e^{-cd^2L^{1/3}}.
\]
\end{theorem}
\begin{proof}
For $|\delta|<1$ define
\[
 n(\delta)=\left(\frac{(1+\delta)^s+(1-\delta)^s}{2}\right)^{1/s},
 \quad w_\delta(x)=\frac{1+\delta x}{n(\delta)},\quad
 m(\delta)=n(\delta)^{-1}.
\]
The one-site matrix is
\[
 H(\delta)=\frac1{n(\delta)}
 \begin{pmatrix}1&\delta\sqrt\rho\\\delta\sqrt\rho&\rho\end{pmatrix},
 \qquad D=\diag(1,\rho).
\]
It is positive definite with top eigenvalue
\[
 \kappa(\delta)=
 \frac{1+\rho+\sqrt{(1-\rho)^2+4\rho\delta^2}}{2n(\delta)}.
\]
For $\tau(\delta)=\|D^{1/2}H(\delta)^{1/2}\|$, the square is the
top eigenvalue of $D^{1/2}H(\delta)D^{1/2}$, so
\[
 \tau(\delta)^2=
 \frac{1+\rho^2+\sqrt{(1-\rho^2)^2+4\rho^2\delta^2}}{2n(\delta)}.
\]
These positive functions have even analytic logarithms near zero.
Taylor expansion with bounded fourth derivative gives, at $s=s_c$,
\begin{align*}
 \log n(\delta)&=\frac{\rho}{1-\rho}\delta^2+O(\delta^4),\\
 \log\kappa(\delta)&=O(\delta^4),\\
 \log\tau(\delta)&=-\frac{\rho}{2(1-\rho^2)}\delta^2+O(\delta^4).
\end{align*}
Indeed, the respective square-root numerators divided by two are
$1+\rho\delta^2/(1-\rho)+O(\delta^4)$ and
$1+\rho^2\delta^2/(1-\rho^2)+O(\delta^4)$.
Consequently there are $0<\delta_*<1$, $C\ge1$ and $c>0$ such that,
for $|\delta|\le\delta_*$,
\[
 |\log\kappa(\delta)|\le C\delta^4,
 \qquad\log\tau(\delta)\le-c\delta^2.
\]
For example, $c=\rho/[4(1-\rho^2)]$ works after reducing $\delta_*$.
Choose a fixed $0<d\le\delta_*$ with $Cd^4\le\log(B-1)$, possible
because $B>2$. Put
\[
 \delta_L=dL^{-1/3},\quad a_L=\kappa(\delta_L)^{-L},\quad
 R_L=\prod_{i=1}^Lw_{\delta_L}(x_i),\quad
 Z_L=1+a_Lm(\delta_L)^L,\quad r_L=\frac{1+a_LR_L}{Z_L}.
\]
The coefficient balances the two top block eigenvalues exactly.
Since $|\log a_L|\le Cd^4L^{-1/3}\le\log(B-1)$,
we have $a_L\le B-1$ and $a_L\to1$.
Product measure gives $\|R_L\|_s=1$ and
$\mathbb ER_L=m(\delta_L)^L$, with $0<m(\delta_L)<1$.
Therefore $\mathbb Er_L=1$, $\|r_L\|_s\le(1+a_L)/Z_L\le B$,
and $r_L\ge1/Z_L\ge1/B$ for every $L\ge1$.

Let $X_L=H(\delta_L)^{\otimes L}$. Spectral equivalence gives
$F_L=Z_LS_{r_L}=P_\rho+a_LX_L$. Both summands have top eigenvalue one.
For the $N$ by $2N$ matrix
$\mathcal C_L=[P_\rho^{1/2},\sqrt{a_L}X_L^{1/2}]$, $N=2^L$,
the Gram matrix $\mathcal C_L^*\mathcal C_L$ has those summands on
its diagonal and off-diagonal block of norm
\[
 \eta_L=\sqrt{a_L}\,\tau(\delta_L)^L
 \le\sqrt{B-1}\,e^{-cd^2L^{1/3}}.
\]
Its block diagonal has two leading eigenvalues exactly one. The
off-diagonal self-adjoint perturbation has norm $\eta_L$, whence
$\lambda_1\le1+\eta_L$ and $\lambda_2\ge1-\eta_L$.
The $N$ positive eigenvalues of $\mathcal C_L\mathcal C_L^*=F_L$
are precisely the positive eigenvalues of $\mathcal C_L^*\mathcal C_L$;
the latter has $N$ additional zeros. Thus the bounds apply to the first
two eigenvalues of $F_L$, including at $L=1$. Since $F_L\ge P_\rho$,
$\lambda_1(F_L)\ge1$ and $\gap(F_L)\le2\eta_L$.
Scalar normalization and $AA^*/A^*A$ preserve the ratio. Strict entrywise
positivity of $T_{r_L}$ gives its strictly positive finite-dimensional gap.
\end{proof}

The uniform component is a spectral channel here, rather than an additive
error to discard. At the critical exponent the tilted top eigenvalue has
zero quadratic variation, while its cross norm with the uniform channel
has negative quadratic variation. The dimension-dependent tilt separates
these orders, and balancing removes the residual diagonal mismatch.
The small-budget estimate uses classical hypercontractivity and elementary
perturbation theory; the large-budget family asserts neither matching
marginals nor fixed-range interactions. No optimality of the sufficient
budget $B>2$ is claimed.
